% ID: SEL-D0011-BA-E02
% ESTADO: cotejada. Segunda pasada visual completa y borrador compilado revisados.
% Fuente: Selectividad/matematicas_aplicadas/EXTRA_SUPLENTE_Modelo-1-Enunciados (Mat-Ap-2021).pdf
% SHA-256: 4a6c7d34ff35c3e7678c52970b4d4558cee2c35224b46c66035cb0bc8d5a7ec8
% Página: 1; bloque A; ejercicio original 2.

Se consideran las matrices $A=\begin{pmatrix}a&4\\6&8\end{pmatrix}$,
$B=\begin{pmatrix}2&2\\3&3\end{pmatrix}$ y $C=\begin{pmatrix}1&2\end{pmatrix}$.
\begin{enumerate}[label=\alph*)]
  \item \textbf{(0.5 puntos)} Calcule el valor del parámetro $a$ para que
  la matriz $A$ no tenga inversa.
  \item \textbf{(1.25 puntos)} Para $a=3$, resuelva la ecuación matricial
  $X\cdot A-X\cdot B=C$.
  \item \textbf{(0.75 puntos)} Para $a=3$, compruebe que $A^2=11\cdot A$
  y exprese $A^8$ en función de la matriz $A$.
\end{enumerate}
